\subsection{The common rank-four space is an actual space}
\label{subsec:sr-four-common-space-construction}

We construct the common space used in
Proposition~\ref{prop:sr-four-core}.
Let \(\Gamma\) be the actual finitely ramified arithmetic
quotient of the Case II tower, with fixed multiplier \(\nu\),
ordered residual characters \(\bar\chi_i\), and distinguished
\(g_0\) at which these characters are pairwise distinct.
All coefficient extensions in this construction are finite.

\paragraph{The formal ring and its two-dimensional grading.}
The universal normalized representation ring \(R\), together
with the four prescribed characters of every ordinary local
group, is complete Noetherian local with finite residue field.
Here is the required finiteness argument.
Every finite-index subgroup of the finitely ramified arithmetic
group has finite-dimensional \(H^1(-,\mathbb F_\ell)\), by the
class-field/Kummer finiteness already proved for these groups.
The maximal pro-\(\ell\) quotient of the residual kernel is
therefore finitely generated, by the pro-\(\ell\) Burnside basis
criterion. Matrix entries for finitely many such generators
and finite residual coset representatives give finitely many
formal variables. Impose their group relations, fixed multiplier,
diagonality at \(g_0\), and killed-inertia relations by a closed
ideal. The local character factors are finite power-series
rings modulo their finite torsion relations, by local class
field theory. Their completed tensor product with the matrix
ring is Noetherian. This constructs \(R\), rather than assuming
a formal parameter ring at each residual point separately.

Put
\[
 v_0=(1,1),\quad v_1=(0,1),\quad
 v_2=(1,0),\quad v_3=(0,0),\qquad w_{ij}=v_i-v_j.
\]
Diagonal symplectic conjugation gives a
\(\mathbb Z^2\)-grading: \(r(g)_{ij}\) has weight \(w_{ij}\)
and every ordinary character has weight zero.
Choose finitely many homogeneous matrix-entry differences and
character differences generating \(\mathfrak m_R\) modulo
\(\mathfrak m_R^2+\mathfrak m_{\mathcal O}R\).
They give a homogeneous presentation
\(R=\mathcal O[[z_1,\ldots,z_b]]/\mathcal I\).
The defining identities are homogeneous, so the closed ideal
\(\mathcal I\) is homogeneous.
Write \(R_w\) for its closed weight-\(w\) part, and set
\[
                 B_{\rm c}=R_0,\qquad
                 A_{\rm c}=\bigoplus_{w\in\mathbb Z^2}R_w
                                      \subset R.          \tag{R4S1}
\]
The direct sum means finite sums of homogeneous elements.

For clarity, the elementary finiteness assertion used here
works for any finite list of integral weights.
The monoid of exponent vectors of weight zero is finitely
generated: intersect the nonnegative orthant with the kernel
of the integral weight matrix, triangulate its rational
polyhedral cone, and reduce coefficients modulo integral
generators of each simplicial cone.
The same argument makes the exponents of any fixed weight
a finitely generated module over this monoid.
Consequently \(B_{\rm c}\) is a quotient of a power-series
ring on finitely many balanced monomials, and is complete
Noetherian local; \(A_{\rm c}\) is generated by the \(z_i\)
over \(B_{\rm c}\); and each \(R_w\) is a finite
\(B_{\rm c}\)-module.
The topology induced on a fixed weight from \(R\) is its
\(\mathfrak m_{B_{\rm c}}\)-adic topology.
Indeed a sufficiently long monomial of that fixed weight
contains arbitrarily many nonconstant balanced factors,
apart from one of finitely many bounded remainders.
These observations prove both the finite-type assertion and
the continuity assertions for the universal matrix entries.

\begin{srlemma}[Balanced rank-four products are trace expressions]
\label{lem:sr-four-balanced-traces}
Let \(r:\Gamma\to\operatorname{GSp}_4(C)\) be a representation
over a commutative ring, with multiplier in a subring \(E\).
Suppose \(r(g_0)\) is diagonal, its four entries and their
pairwise inverse differences belong to \(E\), and
\(\operatorname{tr}r(g)\in E\) for every \(g\).
Then every product
\[
                \prod_{a=1}^m r(g_a)_{i_a j_a},
       \qquad \sum_a(v_{i_a}-v_{j_a})=0 ,
\]
belongs to \(E\).
\end{srlemma}
\begin{proof}
The spectral projectors \(e_i=E_{ii}\) are polynomials in
\(r(g_0)\) over \(E\).
The \(E\)-algebra spanned by the group matrices therefore
contains them. Every directed-cycle product is a trace
\[
 \operatorname{tr}
   (e_{i_1}r(g_1)e_{i_2}r(g_2)\cdots e_{i_m}r(g_m)),
\]
and belongs to \(E\), since this algebra consists of
\(E\)-linear combinations of group matrices.

For a general balanced product draw its directed edges
\(i_a\to j_a\), and let \(b_i\) be the number of outgoing
minus incoming edges at vertex \(i\).
Balance and \(\sum_i b_i=0\) give
\[
                         b=t(1,-1,-1,1)
                                                        \tag{R4S2}
\]
for an integer \(t\).
Symplectic inversion replaces an entry \(i\to j\) by the
entry \(3-j\to3-i\) of \(r(g^{-1})\), multiplied by its
known signed multiplier. Its weight is unchanged.
For a crossing edge from \(\{0,3\}\) to \(\{1,2\}\),
this replacement subtracts \((1,-1,-1,1)\) from \(b\);
for the reverse crossing it adds this vector.
If \(t>0\), the number of outward crossings minus inward
crossings is \(2t\), so at least \(t\) outward crossings
can be replaced. If \(t<0\), use \(-t\) inward crossings.
After these replacements every vertex divergence is zero.
The resulting directed multigraph is a disjoint union
of directed cycles.
Their products belong to \(E\) by the first paragraph.
All introduced multipliers are in \(E\), proving the claim.
\end{proof}

The ring \(B_{\rm c}\) is topologically generated by the
characteristic coefficients and the ordinary character
values. To verify this precise generation assertion, take
their closed subring \(C\) in \(B_{\rm c}\).
The characteristic polynomial of \(g_0\) has simple prescribed
residual roots. Hensel iteration gives its four ordered
roots and their inverse differences in \(C\).
This iteration is legitimate in the induced topology:
\(C\) is closed in the complete ring \(B_{\rm c}\), and a
unit with nonzero constant residue has its inverse in \(C\)
by its convergent geometric series.
The roots are the actual diagonal entries by uniqueness.
Lemma~\ref{lem:sr-four-balanced-traces} puts every balanced
matrix-entry monomial in \(C\).
Expanding diagonal-entry differences adds only constant terms;
the character differences have weight zero and already
belong to \(C\).
Thus every balanced monomial in the chosen \(z_i\) belongs
to \(C\). The preceding monoid description of \(R_0\)
gives \(C=B_{\rm c}\).

\paragraph{The common space receives the nonsplit residual points.}
This assertion requires more than the universal property of \(R\)
at its diagonal vertex. We give the extension argument.

Let \(C\) first be an Artinian local coefficient algebra,
\(\mathfrak m_C^M=0\), and let a normalized representation
\(q\) have diagonal residual entries \(\bar\chi_i\).
Suppose a linear functional
\(\tau:\mathbb Z^2\to\mathbb Z\) makes
\[
 q(g)_{ij}\in\mathfrak m_C
       \quad\text{whenever }i\ne j,\ \tau(w_{ij})\le0.
                                                        \tag{R4S3}
\]
This includes exactly the two allowed residual support cones:
one may take \(\tau(x,y)=-2x-y\) or \(-x-2y\),
respectively.
Take \(c\ge1\) bounding the negatives of the finitely many
\(\tau(w_{ij})\), and use the subring of \(C((t))\)
\[
 S=\sum_{j=0}^{M-1}
             \mathfrak m_C^{\,j}t^{-cj}C[[t]].             \tag{R4S4}
\]
It is a ring because \(\mathfrak m_C^M=0\).
It is generated over \(C[[t]]\) by finitely many
\(\mathfrak m_Ct^{-c}\) generators, and is finite as a
\(C[[t]]\)-module. Thus it is Noetherian.
It is complete local: its maximal ideal is generated by
\(t\) and \(\mathfrak m_Ct^{-c}\), and the latter ideal
is nilpotent, so this topology is equivalent to the
\(t\)-adic topology. Its residue field is that of \(C\).

Scale \(q(g)_{ij}\) by \(t^{\tau(w_{ij})}\).
Equation (R4S3) makes every entry an element of \(S\),
and every off-diagonal entry belongs to its maximal ideal.
Local character values are left unchanged.
The scaling respects matrix multiplication and the fixed
multiplier, since
\(v_i+v_{3-i}=(1,1)\).
The scaled representation therefore gives a map \(R\to S\).

Under this map a weight-\(w\) expression has the form
\(t^{\tau(w)}a\), \(a\in C\).
This also holds for completed homogeneous expressions:
their polynomial approximations converge in \(S\), while
a coefficient at one fixed Laurent exponent is discrete,
so it eventually stabilizes.
Remove \(t^{\tau(w)}\) on each homogeneous part.
The resulting maps combine, by the grading, into
\[
                          A_{\rm c}\longrightarrow C
                                                        \tag{R4S5}
\]
with the original matrices \(q\) and ordinary characters.
For a complete local \(C\), apply this construction modulo
every power of \(\mathfrak m_C\).
Its values on the chosen generators are forced, so the maps
are compatible; completeness gives (R4S5).
Uniqueness follows from the generation of \(B_{\rm c}\)
just proved and the finite-type generation of \(A_{\rm c}\).
This proves the required common-space property for the
actual nonsplit residual supports.

\paragraph{Adjoin the actual flags.}
For each ordinary place adjoin the four coordinates
\(a_v,b_v,c_v,d_v\) of
\[
 n_v=
 \begin{pmatrix}
 1&0&0&0\\a_v&1&0&0\\b_v&d_v&1&0\\
 c_v&b_v-d_va_v&-a_v&1
 \end{pmatrix}.
\]
Impose that \(n_v^\#r(g)n_v\) is upper triangular with
the specified diagonal characters on the local group,
and that these characters are \(1\) on the specified
\(\ell\)-primary torsion subgroup.
Let \(A\) be this quotient of
\(A_{\rm c}[\mathbf a,\mathbf b,\mathbf c,\mathbf d]\),
and replace \(B_{\rm c}\) by its image \(B\) in \(A\).
This is a quotient by an ideal in the Noetherian complete
ring \(B_{\rm c}\), so \(B\) remains complete local.
The split residual representation and its coordinate flags
give a point; hence \(A\ne0\).
By construction \(B\hookrightarrow A\), \(A\) is of finite
type over \(B\), its representation and flag entries
generate it, and its characteristic coefficients and
ordinary characters belong to \(B\) and topologically
generate \(B\).
Every allowed normalized lift with a residual point in
either support cone has its unique map from \(A\), by
(R4S5) and the flag equations.

Completion at a residual point therefore represents its
actual normalized deformation problem.
The two inverse transformations are evaluation of the
universal matrices/flags and the map just constructed.
This proves the common-space and completion assertions
retained in Section~\ref{subsec:sr-rank-four-core}.

The dimension transport used there is the usual finite-type
dimension formula over the complete Noetherian base.
For a prime \(\mathfrak p\) beneath a closed residual point,
apply that formula to the domain \(A/\mathfrak p\)
over \(B/(\mathfrak p\cap B)\); the residual field extension
is finite. It gives the sum of the base dimension and the
fraction-field transcendence degree.
Localization and completion preserve these component
dimensions, since complete Noetherian coefficient rings are
excellent and their formal local fibers are equidimensional.
Thus a bound on the number of additional field generators
is precisely the completed-quotient bound used in
Lemma~\ref{lem:sr-four-trace-fiber}.
No smoothness of a reducible representation is needed.

Finally a polynomial residual path gives a map
\(A\to k[t]\): apply the preceding property first to
\(k[[t]]\). The characteristic coefficients and ordinary
characters on a residual path are their constant prescribed
values, so its restriction to \(B\) is the residue map.
The remaining generators are the displayed polynomial
matrix and flag entries; the image is therefore in
\(k[t]\). This proves the actual line property used to
transport components along residual paths.

The frozen-source construction expanded here is
\srcid{AdicExtensionMatrices.SxUnivD},
\srcid{GenZero}, \srcid{NearArt}, \srcid{QuotPowNil},
\srcid{PreOfCore}, \srcid{PreImage}, and \srcid{PreLine},
assembled in
\srcid{WeakShimuraTheory.CompactIntegralFields.SpaceEx}.
The graph proof of the trace assertion expands
\srcid{MonoZero} and \srcid{GdotSelf}.
